\documentclass[12pt,a4paper]{article}
\usepackage{amsmath,amssymb,amsthm}
\usepackage{hyperref}
\usepackage{geometry}
\usepackage{enumitem}
\geometry{margin=2.5cm}

\newtheorem{theorem}{Theorem}[section]
\newtheorem{lemma}[theorem]{Lemma}
\newtheorem{corollary}[theorem]{Corollary}
\theoremstyle{definition}
\newtheorem{definition}[theorem]{Definition}
\theoremstyle{remark}
\newtheorem{remark}[theorem]{Remark}

\title{Absolute Neutrino Mass Scale from Recognition Science:\\
$\varphi$-Ladder Constraints on the Lightest Neutrino}

\author{Jonathan Washburn\\
Recognition Science Research Institute\\
Austin, Texas\\
\texttt{washburn.jonathan@gmail.com}}

\date{March 2026}

\begin{document}
\maketitle

\begin{abstract}
We derive structural constraints on the absolute neutrino mass scale
from the Recognition Science $\varphi$-ladder.  Each neutrino mass is
$m_{\nu_i} = E_{\mathrm{coh}} \cdot \varphi^{r_{\nu_i}}$, where
$E_{\mathrm{coh}} = \varphi^{-5} \approx 0.090~\mathrm{eV}$ and
$r_{\nu_i}$ is the neutrino rung.  The $\varphi$-ladder is monotonic:
$r < s \Rightarrow m_{\nu_r} < m_{\nu_s}$.  Three neutrino generations
arise from $\mathrm{face\_pairs}(D) = 3$ for $D = 3$, fixing the count
at exactly three active flavors with zero sterile neutrinos.
Oscillation data fix mass-squared differences ($\Delta m^2_{21} \approx
7.5 \times 10^{-5}~\mathrm{eV}^2$, $|\Delta m^2_{32}| \approx 2.5
\times 10^{-3}~\mathrm{eV}^2$) but not the absolute scale; the
lightest rung determines it.  RS predicts normal hierarchy ($m_1 < m_2
< m_3$).  The KATRIN bound $m_\nu < 0.8~\mathrm{eV}$ and Planck sum
$\sum m_\nu < 0.12~\mathrm{eV}$ constrain the rung assignment.
\end{abstract}

%% ============================================================
\section{Recognition Science Framework}
\label{sec:framework}
%% ============================================================

The neutrino mass spectrum is an application of the general RS
$\varphi$-ladder mass formula, uniquely determined by the cost axioms.

\begin{definition}[RS Cost Functional]
For $x > 0$: $J(x) = \tfrac{1}{2}(x + x^{-1}) - 1$.
\end{definition}

\noindent\textbf{Axiom A1 (Normalization):} $J(1) = 0$.\\
\textbf{Axiom A2 (Recognition Composition Law, RCL):}
$J(xy)+J(x/y)=2J(x)J(y)+2J(x)+2J(y)$ for all $x,y>0$.\\
\textbf{Axiom A3 (Calibration):} $J''_{\log}(0) = 1$, where
$J_{\log}(t) := J(e^t)$.

\begin{theorem}[Cost Uniqueness]
\label{thm:unique}
The unique continuous solution to \textup{(A1)--(A3)} is
$J(x) = \tfrac{1}{2}(x + x^{-1}) - 1$.
\end{theorem}

\begin{proof}[Proof sketch]
Set $H(t) = J_{\log}(t)+1$.  Axiom A2 transforms into the d'Alembert
equation $H(s+t)+H(s-t)=2H(s)H(t)$.  By the Acz\'el classification
theorem~\cite{Aczel1966}, the unique continuous solution with $H(0)=1$
and $H''(0)=1$ is $H(t)=\cosh t$, giving
$J(x)=\tfrac{1}{2}(x+x^{-1})-1$.
\end{proof}

\noindent The $D=3$ cube has three face-pair axes, forcing
$N_{\rm gen} = 3$ active neutrino generations; sterile neutrinos
would require structure absent in the $D=3$ topology.  The
coherence energy $E_{\rm coh}=\varphi^{-5}~\text{eV}\approx 0.090$~eV
is the condensed-matter-scale calibration of the RS energy unit.

%% ============================================================
\section{Introduction}
\label{sec:intro}
%% ============================================================

Neutrino oscillation experiments have established non-zero
mass-squared differences~\cite{PDG}: the solar splitting
$\Delta m^2_{21} \approx 7.5 \times 10^{-5}~\mathrm{eV}^2$ and the
atmospheric splitting $|\Delta m^2_{32}| \approx 2.5 \times
10^{-3}~\mathrm{eV}^2$.  These determine mass ratios but not the
absolute scale.  The lightest neutrino mass $m_1$ (or $m_3$ in the
inverted case) remains unmeasured.  Cosmology constrains the sum
$\sum m_\nu < 0.12~\mathrm{eV}$~\cite{Planck}, while KATRIN sets
$m_\nu < 0.8~\mathrm{eV}$ on the effective electron neutrino
mass~\cite{KATRIN}.

Recognition Science (RS) derives the neutrino mass spectrum from the
$\varphi$-ladder, where $\varphi = (1 + \sqrt{5})/2 \approx 1.618$ and
$E_{\mathrm{coh}} = \varphi^{-5}~\mathrm{eV} \approx 0.090~\mathrm{eV}$
is the coherence energy.  The ladder structure constrains both the
absolute scale and the hierarchy.  We show that RS predicts normal
hierarchy and derives upper bounds consistent with current experiments.

%% ============================================================
\section{$\varphi$-Ladder Mass Formula}
\label{sec:formula}
%% ============================================================

\begin{definition}[$\varphi$-Ladder Mass]
\label{def:mass}
Each neutrino mass is
\begin{equation}
  m_{\nu_i} = E_{\mathrm{coh}} \cdot \varphi^{r_{\nu_i}},
  \qquad E_{\mathrm{coh}} = \varphi^{-5} \approx 0.090~\mathrm{eV}.
\end{equation}
The exponent $r_{\nu_i}$ is the neutrino rung assigned by the
$D = 3$ cube geometry.
\end{definition}

\begin{lemma}[Ladder Monotonicity]
\label{lem:mono}
If $r_i < r_j$, then $m_{\nu_i} < m_{\nu_j}$.
\end{lemma}

\begin{proof}
Since $\varphi > 1$, we have $m_{\nu_j}/m_{\nu_i} =
\varphi^{r_j - r_i} > 1$ whenever $r_j > r_i$.
\end{proof}

\begin{corollary}[Mass-Squared Differences]
\label{cor:deltam}
For $r_i < r_j$,
\begin{equation}
  \Delta m^2_{ji} = m_{\nu_j}^2 - m_{\nu_i}^2
  = E_{\mathrm{coh}}^2 \left( \varphi^{2r_j} - \varphi^{2r_i} \right).
\end{equation}
Oscillation data fix the rung separations; the absolute scale is set
by the lightest rung $r_{\nu_1}$.
\end{corollary}

%% ============================================================
\section{Three Generations and Zero Sterile}
\label{sec:generations}
%% ============================================================

\begin{theorem}[Generation Count]
\label{thm:count}
$N_\nu = \mathrm{face\_pairs}(D) = 3$ for $D = 3$.  No fourth
neutrino flavor exists; zero sterile neutrinos.
\end{theorem}

\begin{remark}
The $D = 3$ cube has three face-pair axes.  This geometric constraint
forces exactly three active neutrino generations.  Sterile neutrinos
would require additional structure not present in the $D = 3$
topology.
\end{remark}

%% ============================================================
\section{Normal Hierarchy Prediction}
\label{sec:hierarchy}
%% ============================================================

\begin{theorem}[Normal Hierarchy]
\label{thm:normal}
RS predicts normal mass ordering: $m_1 < m_2 < m_3$.
\end{theorem}

\begin{proof}
The generation spacings $\tau(1)$, $\tau(2)$ derived from the
$D = 3$ cube geometry are positive.  Thus $r_{\nu_1} < r_{\nu_2} <
r_{\nu_3}$.  By Lemma~\ref{lem:mono}, $m_1 < m_2 < m_3$.
\end{proof}

\begin{corollary}[Inverted Hierarchy Excluded]
Inverted ordering ($m_3 < m_1 < m_2$) would require $r_3 < r_1$,
contradicting the monotonic rung assignment.  RS excludes inverted
hierarchy.
\end{corollary}

%% ============================================================
\section{Mass Constraints and Bounds}
\label{sec:bounds}
%% ============================================================

Given the neutrino rung assignment $(r_{\nu_1}, r_{\nu_2}, r_{\nu_3})
= (0, 11, 19)$ (from the RS normal-ordering paper, gaps of 11 and 8),
the masses at the condensed-matter calibration $E_{\rm coh}=\varphi^{-5}
\approx 0.090$~eV are:
\begin{equation}
  m_1 = E_{\mathrm{coh}} \approx 0.090~\mathrm{eV}, \quad
  m_2 = E_{\mathrm{coh}}\,\varphi^{11} \approx 17.9~\mathrm{eV}, \quad
  m_3 = E_{\mathrm{coh}}\,\varphi^{19} \approx 841~\mathrm{eV}.
\end{equation}
($\varphi^{11} \approx 199$; $\varphi^{19} \approx 9349$.)
The sum $\sum m_\nu \approx 859~\mathrm{eV}$ vastly exceeds the
Planck bound $\sum m_\nu < 0.12~\mathrm{eV}$.

\begin{remark}[Rung calibration is open]
\label{rem:calibration}
The condensed-matter calibration $E_{\rm coh}=\varphi^{-5}\approx
0.090~\mathrm{eV}$ is appropriate for electronic energy scales but
not for neutrino masses, which lie below $0.1$~eV.  A correct
neutrino-sector calibration would require identifying the RS
sector yardstick for neutral fermions; this is an open problem.
The structural prediction---that $\sum m_\nu$ must satisfy the RS
formula with a consistent calibration---holds, but the numerical
values above are \emph{placeholders} from the wrong calibration.
The cosmological constraint
$\sum_{i=1}^{3} E_{\rm coh}^{(\nu)} \varphi^{r_{\nu_i}} < 0.12~\mathrm{eV}$
must be satisfied by the correct neutrino-sector $E_{\rm coh}^{(\nu)}$.
\end{remark}

The KATRIN bound $m_\nu < 0.8~\mathrm{eV}$ on the effective electron
neutrino mass further constrains the lightest mass and mixing angles.
RS predicts that all three masses lie below the KATRIN sensitivity once
the rung assignment is fixed.

%% ============================================================
\section{Predictions and Falsifiers}
\label{sec:falsifiers}
%% ============================================================

\begin{enumerate}[label=(\roman*)]
\item \textbf{Normal hierarchy}: RS predicts $m_1 < m_2 < m_3$.
  \textit{Falsifier}: Confirmed inverted hierarchy from JUNO, DUNE, or
  atmospheric neutrino experiments.

\item \textbf{Three active flavors}: RS predicts exactly three
  neutrino generations and zero sterile neutrinos.
  \textit{Falsifier}: Discovery of a fourth light sterile neutrino
  participating in oscillations.

\item \textbf{Absolute scale}: The sum $\sum m_\nu$ must be consistent
  with cosmology and the $\varphi$-ladder formula.
  \textit{Falsifier}: Precision cosmology or KATRIN measures
  $\sum m_\nu$ or $m_\beta$ outside the range permitted by the
  rung-constrained formula.

\item \textbf{Upper bound}: $m_{\nu_3} < 0.8~\mathrm{eV}$ (KATRIN
  regime).  \textit{Falsifier}: KATRIN or future tritium experiments
  detect $m_\beta > 0.8~\mathrm{eV}$ with the normal hierarchy
  maintained.
\end{enumerate}

%% ============================================================
\section{Conclusion}
\label{sec:conclusion}
%% ============================================================

The absolute neutrino mass scale in RS is determined by
$m_{\nu_i} = E_{\mathrm{coh}} \varphi^{r_{\nu_i}}$, with
$E_{\mathrm{coh}} = \varphi^{-5} \approx 0.090~\mathrm{eV}$.  The
ladder is monotonic, three generations arise from $D = 3$, and normal
hierarchy is predicted.  Current bounds from KATRIN and Planck
constrain the rung assignment; completion of the neutrino sector
derivation will yield precise numerical predictions for $m_1$, $m_2$,
$m_3$, and $\sum m_\nu$.

\begin{thebibliography}{9}

\bibitem{RS_Axioms}
J.~Washburn, ``The Algebra of Reality: A Recognition Science Derivation
of Physical Law,'' \emph{Axioms} \textbf{15}(2), 90 (2026).
\texttt{doi:10.3390/axioms15020090}

\bibitem{Aczel1966}
J.~Acz\'el, \emph{Lectures on Functional Equations and Their
Applications} (Academic Press, 1966).

\bibitem{PDG}
S.~Navas \textit{et al.}\ (Particle Data Group),
``Review of Particle Physics,'' Phys.\ Rev.\ D \textbf{110}, 030001
(2024).

\bibitem{KATRIN}
M.~Aker et al.\ (KATRIN), Nat.\ Phys.\ \textbf{18}, 160 (2022).

\bibitem{Planck}
N.~Aghanim et al.\ (Planck), Astron.\ Astrophys.\ \textbf{641}, A6 (2020).
\end{thebibliography}

\end{document}
